\documentclass[12pt]{article}
\usepackage[utf8]{inputenc}
\usepackage[T1]{fontenc}
\usepackage[margin=1in]{geometry}
\usepackage{setspace}
\usepackage{xcolor}
\usepackage[colorlinks=true, urlcolor=blue!60!black, linkcolor=blue!60!black]{hyperref}
\setlength{\parindent}{2em}
\setlength{\parskip}{1.25em}
\renewcommand{\baselinestretch}{1.0}
\begin{document}
\begin{center}
{\LARGE \textit{Geduldig v. Aiello}, 417 U.S. 484 (1974)\par}
\vspace{0.8em}
{\large Case Brief\par}
\end{center}
\vspace{1em}

\section{Memory Jogger}

Excluding normal pregnancy from a state-run employee disability insurance program does not constitute sex-based discrimination under the Equal Protection Clause because the program divides potential recipients into ``pregnant women and nonpregnant persons'' rather than men and women, triggering only rational basis review (\href{https://www.courtlistener.com/opinion/109065/geduldig-v-aiello/}{CourtListener}; \href{https://supreme.justia.com/cases/federal/us/417/484/}{Justia}).

\section{Detailed Case Facts}

Since 1946, California operated a mandatory, entirely self-supporting disability insurance system for private-sector employees temporarily unable to work due to mental or physical illness or injury not covered by workers' compensation (\href{https://supreme.justia.com/cases/federal/us/417/484/}{Justia}). Participating employees funded the Unemployment Compensation Disability Fund by contributing 1\% of their annual wages up to a maximum annual contribution of \$85 (raised to \$90 for 1974 and thereafter), with no general state tax subsidies (\href{https://supreme.justia.com/cases/federal/us/417/484/}{Justia}). Eligible employees who had contributed 1\% of at least \$300 in base-period wages received weekly benefits ranging from \$25 to \$105 (increased to \$119 effective July 1, 1974) for up to 26 weeks, beginning on the eighth day of disability or the first day of hospitalization, plus \$12 per day in hospital benefits for up to 20 days (\href{https://supreme.justia.com/cases/federal/us/417/484/}{Justia}).

The program did not insure every disabling condition: it excluded disabilities lasting fewer than eight days without hospitalization, disabilities extending beyond 26 weeks, and disabilities resulting from court commitment as a dipsomaniac, drug addict, or sexual psychopath (\href{https://supreme.justia.com/cases/federal/us/417/484/}{Justia}). Crucially, California Unemployment Insurance Code \S~2626 excluded ``any injury or illness caused by or arising in connection with pregnancy up to the termination of such pregnancy and for a period of 28 days thereafter'' (\href{https://supreme.justia.com/cases/federal/us/417/484/}{Justia}).

Four otherwise-eligible California working women who became disabled due to pregnancy were denied benefits under \S~2626 and sued Dwight Geduldig, Director of the California Department of Human Resources Development (\href{https://supreme.justia.com/cases/federal/us/417/484/}{Justia}). Three of the women suffered abnormal pregnancy complications---Carolyn Aiello (ectopic pregnancy requiring surgery), Elizabeth Johnson (tubal pregnancy requiring surgery), and Augustina Armendariz (miscarriage)---while the fourth, Jacqueline Jaramillo, experienced a normal pregnancy and childbirth (\href{https://supreme.justia.com/cases/federal/us/417/484/}{Justia}).

Ten days before the federal District Court ruled, the California Court of Appeal decided \textit{Rentzer v. Unemployment Insurance Appeals Board}, 32 Cal.\ App.\ 3d 604 (1973), construing \S~2626 to exclude only disabilities accompanying \textit{normal} pregnancy and childbirth (\href{https://supreme.justia.com/cases/federal/us/417/484/}{Justia}). California acquiesced in \textit{Rentzer}, issued revised administrative guidelines, and amended \S~2626 (adding \S~2626.2 effective January 1, 1974) to cover abnormal pregnancy complications (\href{https://supreme.justia.com/cases/federal/us/417/484/}{Justia}). As a result, Aiello, Armendariz, and Johnson were paid their benefits, leaving only Jaramillo's normal-pregnancy claim as a live controversy (\href{https://supreme.justia.com/cases/federal/us/417/484/}{Justia}).

\section{Procedural History}

\begin{enumerate}
  \item \textbf{Initial Actions and Consolidation:} Carolyn Aiello filed a class action in the U.S.\ District Court for the Northern District of California challenging \S~2626 under the Equal Protection Clause of the Fourteenth Amendment (\href{https://supreme.justia.com/cases/federal/us/417/484/}{Justia}). Armendariz, Johnson, and Jaramillo jointly petitioned for a writ of mandate in the California Supreme Court; the State removed that action to federal court under 28 U.S.C.\ \S~1441(b), where the two suits were consolidated (\href{https://supreme.justia.com/cases/federal/us/417/484/}{Justia}).
  \item \textbf{Three-Judge U.S.\ District Court (N.D.\ Cal.):} Because plaintiffs sought to enjoin enforcement of a state statute, a three-judge court was convened under 28 U.S.C.\ \S\S~2281 and 2284 (\href{https://supreme.justia.com/cases/federal/us/417/484/}{Justia}). In \textit{Aiello v. Hansen}, 359 F. Supp. 792 (N.D. Cal. 1973), the divided panel granted summary judgment for plaintiffs, holding that the pregnancy exclusion lacked a rational and substantial relationship to a legitimate state purpose and violated the Equal Protection Clause (\href{https://supreme.justia.com/cases/federal/us/417/484/}{Justia}). The court denied reconsideration in light of \textit{Rentzer} and denied a stay pending appeal (\href{https://supreme.justia.com/cases/federal/us/417/484/}{Justia}).
  \item \textbf{U.S.\ Supreme Court:} On direct appeal, the Supreme Court granted a stay (414 U.S. 897 (1973)) and noted probable jurisdiction (414 U.S. 1110 (1973)) (\href{https://supreme.justia.com/cases/federal/us/417/484/}{Justia}). Following oral argument on March 26, 1974, the Court vacated the stay and reversed on June 17, 1974 (\href{https://supreme.justia.com/cases/federal/us/417/484/}{Justia}; \href{https://www.oyez.org/cases/1973/73-640}{Oyez}).
\end{enumerate}

\section{Judicial Votes}

\textbf{Vote Breakdown:} 6--3, reversing the three-judge District Court (\href{https://supreme.justia.com/cases/federal/us/417/484/}{Justia}; \href{https://tile.loc.gov/storage-services/service/ll/usrep/usrep417/usrep417484/usrep417484.pdf}{Library of Congress U.S.\ Reports}).

\begin{itemize}
  \item \textbf{Majority (6):} Justice Potter Stewart (author), joined by Chief Justice Warren E. Burger, and Justices Byron R. White, Harry A. Blackmun, Lewis F. Powell, Jr., and William H. Rehnquist.
  \item \textbf{Concurrences (0):} None.
  \item \textbf{Dissent (3):} Justice William J. Brennan, Jr. (author), joined by Justices William O. Douglas and Thurgood Marshall.
\end{itemize}

\section{Holding}

\textbf{Judgment:} Stay vacated; judgment of the District Court reversed (\href{https://supreme.justia.com/cases/federal/us/417/484/}{Justia}).

\textbf{Rule Announced:} First, the claims of the three appellees with abnormal pregnancies were mooted by \textit{Rentzer} and the State's payment of their claims (\href{https://supreme.justia.com/cases/federal/us/417/484/}{Justia}). Second, a state's exclusion of normal pregnancy from a comprehensive, self-supporting disability insurance program is not a sex-based classification triggering heightened scrutiny; absent a showing that a pregnancy-based distinction is a mere pretext for invidious sex discrimination, it is evaluated under rational basis review and is constitutional so long as it rationally furthers legitimate state fiscal and administrative interests (\href{https://supreme.justia.com/cases/federal/us/417/484/}{Justia}).

\textbf{Holding vs.\ Dicta:} The precise holding is that California's exclusion of normal pregnancy-related disabilities from its employee-funded disability insurance program did not violate the Fourteenth Amendment's Equal Protection Clause (\href{https://supreme.justia.com/cases/federal/us/417/484/}{Justia}). However, the Court's sweeping conceptual rationale in Footnote 20---that legislative classifications concerning pregnancy do not constitute sex-based classifications because they divide people into ``pregnant women and nonpregnant persons'' unless shown to be a ``mere pretext'' for invidious discrimination---became the decision's most consequential doctrinal rule, later cited in \textit{General Electric Co. v. Gilbert}, 429 U.S. 125 (1976), \textit{Dobbs v. Jackson Women's Health Organization}, 597 U.S. 215 (2022), and \textit{United States v. Skrmetti}, 605 U.S. \_\_\_ (2025) (\href{https://www.supremecourt.gov/opinions/21pdf/19-1392_6j37.pdf}{Dobbs Slip Op.}; \href{https://www.law.cornell.edu/supremecourt/text/23-477}{Cornell LII, Skrmetti}). Conversely, the majority's observation in Footnote 21 that women contributed roughly 28\% of the fund while receiving 38\% of total benefits functioned as supporting empirical context rather than a necessary element of the holding (\href{https://supreme.justia.com/cases/federal/us/417/484/}{Justia}).

\section{Analysis of Opinions}

\subsection{Majority Opinion (Justice Stewart)}

Justice Stewart treated California's disability system as a classic social-welfare insurance program subject to deferential rational basis review under \textit{Williamson v. Lee Optical Co.}, 348 U.S. 483 (1955), \textit{Dandridge v. Williams}, 397 U.S. 471 (1970), and \textit{Jefferson v. Hackney}, 406 U.S. 535 (1972) (\href{https://supreme.justia.com/cases/federal/us/417/484/}{Justia}). Because the Disability Fund had paid out between 90\% and 103\% of its annual revenues in recent years, adding normal pregnancy coverage---which the State estimated would cost \$120.2 million to \$131 million annually (a 33\% to 36\% increase)---would force California to sacrifice its self-supporting structure, lower benefits for covered disabilities, or raise the 1\% contribution rate in a way that would regressively burden low-income workers (\href{https://supreme.justia.com/cases/federal/us/417/484/}{Justia}). Under \textit{Dandridge}, the Equal Protection Clause does not force a state to ``choose between attacking every aspect of a problem or not attacking the problem at all''; excluding normal pregnancy was no different in principle from excluding short-term disabilities under eight days or long-term disabilities beyond 26 weeks (\href{https://supreme.justia.com/cases/federal/us/417/484/}{Justia}).

To avoid heightened sex-discrimination scrutiny under \textit{Reed v. Reed}, 404 U.S. 71 (1971), and \textit{Frontiero v. Richardson}, 411 U.S. 677 (1973), Justice Stewart asserted that the program covered the same package of risks for both sexes: ``There is no risk from which men are protected and women are not. Likewise, there is no risk from which women are protected and men are not'' (\href{https://supreme.justia.com/cases/federal/us/417/484/}{Justia}). In the famous Footnote 20, he reasoned that while only women can become pregnant, normal pregnancy is ``an objectively identifiable physical condition with unique characteristics,'' and the program divided potential recipients into ``pregnant women and nonpregnant persons''---the first exclusively female, the second containing both men and women (\href{https://supreme.justia.com/cases/federal/us/417/484/}{Justia}). Absent proof of invidious pretext, excluding pregnancy was not a sex classification (\href{https://supreme.justia.com/cases/federal/us/417/484/}{Justia}).

\subsection{Dissenting Opinion (Justice Brennan, Joined by Justices Douglas and Marshall)}

Justice Brennan argued that California's statute imposed a stark gender-based double standard that warranted strict judicial scrutiny under his \textit{Frontiero} plurality view (\href{https://supreme.justia.com/cases/federal/us/417/484/}{Justia}). He emphasized that California's disability fund was a pooled-risk humanitarian program that compensated virtually every other disabling condition regardless of cost, voluntariness, or ``normalcy''---including heart attacks, cosmetic surgery, sterilization, wisdom-tooth extraction, and sex- or race-linked conditions such as prostatectomies, circumcision, hemophilia, gout, and sickle-cell anemia (\href{https://supreme.justia.com/cases/federal/us/417/484/}{Justia}). Yet normal pregnancy---a medically disabling condition unique to women that carries the same wage-loss and medical-expense consequences---was singled out for exclusion (\href{https://supreme.justia.com/cases/federal/us/417/484/}{Justia}).

Citing the Equal Employment Opportunity Commission's Title VII guidelines (29 C.F.R.\ \S~1604.10(b)), which treated pregnancy disabilities identically to other temporary disabilities, Justice Brennan maintained that dissimilar treatment based on physical characteristics inextricably linked to one sex ``inevitably constitutes sex discrimination'' (\href{https://supreme.justia.com/cases/federal/us/417/484/}{Justia}). Under heightened scrutiny, saving money cannot justify an invidious classification (\textit{Shapiro v. Thompson}, 394 U.S. 618, 633 (1969)), especially when California could preserve solvency through sex-neutral adjustments---such as raising the contribution rate by roughly 0.364\% and increasing the annual ceiling to \$119 (\href{https://supreme.justia.com/cases/federal/us/417/484/}{Justia}).

\section{Examples: Future Applications}

\begin{enumerate}
  \item \textbf{Same-Side (Rational Basis Applies; Upheld):} A state Medicaid program excludes elective, non-therapeutic abortion funding while funding childbirth and other medical procedures, and challengers argue this discriminates on the basis of sex under the Equal Protection Clause. Relying on \textit{Geduldig} (and later \textit{Dobbs v. Jackson Women's Health Organization}, 597 U.S. 215 (2022)), a court would hold that regulating or declining to fund a medical procedure linked to pregnancy does not trigger heightened sex scrutiny absent a showing of pretext for invidious sex discrimination (\href{https://www.supremecourt.gov/opinions/21pdf/19-1392_6j37.pdf}{Dobbs Slip Op.}; \href{https://supreme.justia.com/cases/federal/us/417/484/}{Justia}).
  \item \textbf{Same-Side (Rational Basis Applies; Upheld):} A state enacts a law prohibiting puberty blockers and hormone therapy for minors when prescribed to treat gender dysphoria, while permitting the same medications for other diagnoses such as precocious puberty. Extending \textit{Geduldig}'s ``lack of identity'' logic, the Supreme Court in \textit{United States v. Skrmetti}, 605 U.S. \_\_\_ (2025), held that although only transgender minors seek the excluded treatments, the law divides minors into those seeking treatment for excluded diagnoses and all other minors (a group including both transgender and nontransgender youth), thus triggering only rational basis review absent invidious pretext (\href{https://www.law.cornell.edu/supremecourt/text/23-477}{Cornell LII, Skrmetti}).
  \item \textbf{Opposite-Side (Heightened Scrutiny or Statutory Liability):} A state agency adopts a mandatory unpaid leave policy requiring pregnant employees to stop working at five months of pregnancy based on assumptions that pregnant women are less reliable or should remain at home. Unlike California's actuarial risk exclusion in \textit{Geduldig}, a policy rooted in archaic sex-role stereotypes about women's capacity to work constitutes pretextual sex discrimination under the Equal Protection Clause (cf.\ \textit{Nevada Department of Human Resources v. Hibbs}, 538 U.S. 721 (2003)) and also violates Title VII as amended by the Pregnancy Discrimination Act of 1978, 42 U.S.C.\ \S~2000e(k) (\href{https://digitalcommons.wcl.american.edu/cgi/viewcontent.cgi?article=1409&context=facsch_bk_contributions}{American U.\ WCL}; \href{https://www.law.georgetown.edu/georgetown-law-journal/wp-content/uploads/sites/26/2020/07/The-Pregnant-Citizen-from-Suffrage-to-the-Present.pdf}{Georgetown L.J.}).
  \item \textbf{Opposite-Side (Heightened Scrutiny; Struck Down):} A state university refuses to hire any applicant who has previously been pregnant, citing a desire to reduce employee turnover, while hiring men who have fathered children. Because the policy overtly penalizes a status exclusively tied to women while exempting similarly situated men, a court would find that the pregnancy classification operates as a transparent pretext for invidious sex discrimination, triggering intermediate scrutiny under \textit{United States v. Virginia}, 518 U.S. 515 (1996), and failing constitutional review (\href{https://supreme.justia.com/cases/federal/us/417/484/}{Justia}; \href{https://www.law.georgetown.edu/georgetown-law-journal/wp-content/uploads/sites/26/2020/07/The-Pregnant-Citizen-from-Suffrage-to-the-Present.pdf}{Georgetown L.J.}).
  \item \textbf{Fence-Sitter (Borderline Case):} A state public-employee health plan excludes coverage for prescription contraceptives while covering prescription drugs for erectile dysfunction and all other preventive care. Under \textit{Geduldig}, the state argues that the plan classifies by medical condition rather than sex (dividing employees into contraceptive users and non-users) to contain costs; plaintiffs counter that covering male-specific reproductive medications while excluding female-specific preventive contraception is a textbook pretext for sex discrimination under \textit{Geduldig}'s own Footnote 20 exception (\href{https://digitalcommons.wcl.american.edu/cgi/viewcontent.cgi?article=1409&context=facsch_bk_contributions}{American U.\ WCL}; \href{https://supreme.justia.com/cases/federal/us/417/484/}{Justia}).
\end{enumerate}

\section{Critique}

\textbf{Originalist and Textualist Perspectives:} Conservative originalist and formalist scholars generally defend \textit{Geduldig}'s reluctance to expand heightened Equal Protection scrutiny beyond facially suspect classifications. From an anti-classification standpoint, the Equal Protection Clause polices formal legal categories rather than disparate actuarial impacts in social-welfare design, leaving fiscal trade-offs in state benefit programs to elected legislatures (\href{https://constitution.heritage.org/essays/amdt-14-s1-e/}{Heritage Guide}; \href{https://harvardlawreview.org/print/vol-138/making-equal-protection-protect/}{Harvard L. Rev.}). Moreover, strict originalists such as Justice Thomas (as noted in his \textit{Skrmetti} concurrence) question whether the Fourteenth Amendment's original public meaning in 1868 mandated heightened scrutiny for sex classifications at all, making judicial deference to legislative lines over biological conditions even more compelling (\href{https://www.theregreview.org/2025/08/04/mayeri-shanor-trans-rights-and-the-future-of-sex-discrimination-law/}{The Regulatory Review}). Yet even within formalist ranks, critics point out that \textit{Geduldig}'s ``pregnant women and nonpregnant persons'' syllogism is analytically strained because a classification defined by a trait biologically exclusive to one sex operates as a perfect proxy for that sex (\href{https://verdict.justia.com/2019/10/22/the-pregnant-pause-placing-elizabeth-warrens-experience-of-pregnancy-discrimination-in-historical-context}{Justia Verdict}).

\textbf{Progressive, Feminist, and Doctrinal Perspectives:} \textit{Geduldig} is among the most widely criticized decisions in modern equal protection jurisprudence (\href{https://digitalcommons.wcl.american.edu/cgi/viewcontent.cgi?article=1409&context=facsch_bk_contributions}{American U.\ WCL}). Feminist and progressive scholars---including Reva Siegel, Neil Siegel, Cary Franklin, and Maya Manian---argue that \textit{Geduldig} enshrined an ``invisible male norm'' in workplace benefits: by treating male-typical disabilities (including prostatectomies) as the neutral baseline and pregnancy as an anomalous ``extra'' cost, the Court masked structural subordination as formal equality (\href{https://digitalcommons.wcl.american.edu/cgi/viewcontent.cgi?article=1409&context=facsch_bk_contributions}{American U.\ WCL}; \href{https://scholarship.law.duke.edu/cgi/viewcontent.cgi?article=1450&context=dlj}{Duke L.J.}; \href{https://journals.library.wustl.edu/lawreview/article/5549/galley/22382/view/}{Washington U.\ L.\ Rev.}). When the Court extended \textit{Geduldig}'s logic to Title VII in \textit{General Electric Co. v. Gilbert}, 429 U.S. 125 (1976), Congress swiftly repudiated it by enacting the Pregnancy Discrimination Act of 1978 (\href{https://digitalcommons.wcl.american.edu/cgi/viewcontent.cgi?article=1409&context=facsch_bk_contributions}{American U.\ WCL}). Scholars further argue that later Equal Protection cases such as \textit{United States v. Virginia} and \textit{Hibbs} undermined \textit{Geduldig}'s premises, making the Roberts Court's revival of \textit{Geduldig} in \textit{Dobbs} and \textit{Skrmetti} a major setback for substantive sex and gender equality (\href{https://www.law.georgetown.edu/georgetown-law-journal/wp-content/uploads/sites/26/2020/07/The-Pregnant-Citizen-from-Suffrage-to-the-Present.pdf}{Georgetown L.J.}; \href{https://law.yale.edu/sites/default/files/documents/faculty/papers/franklin-siegel-equality-emerges-chapter-2-endnotes.pdf}{Yale Law}; \href{https://www.theregreview.org/2025/08/04/mayeri-shanor-trans-rights-and-the-future-of-sex-discrimination-law/}{The Regulatory Review}).

\textbf{Underlying Competing Values:} \textit{Geduldig} reflects a deep collision between \textbf{formal anti-classification and legislative fiscal deference} (states must have latitude to design affordable, incremental social insurance programs without courts mandating comprehensive coverage) and \textbf{substantive anti-subordination} (equal citizenship requires dismantling workplace and benefit structures that systematically penalize women for biological reproductive capacity).

\section{Key Quotations}

\begin{quote}
``There is no evidence in the record that the selection of the risks insured by the program worked to discriminate against any definable group or class in terms of the aggregate risk protection derived by that group or class from the program. There is no risk from which men are protected and women are not. Likewise, there is no risk from which women are protected and men are not.''
\end{quote}
\noindent --- Justice Potter Stewart, Opinion of the Court, 417 U.S. at 496--497 (\href{https://supreme.justia.com/cases/federal/us/417/484/}{Justia}; \href{https://www.courtlistener.com/opinion/109065/geduldig-v-aiello/}{CourtListener}).

\begin{quote}
``While it is true that only women can become pregnant, it does not follow that every legislative classification concerning pregnancy is a sex-based classification like those considered in \textit{Reed}, \textit{supra}, and \textit{Frontiero}, \textit{supra}. Normal pregnancy is an objectively identifiable physical condition with unique characteristics. Absent a showing that distinctions involving pregnancy are mere pretexts designed to effect an invidious discrimination against the members of one sex or the other, lawmakers are constitutionally free to include or exclude pregnancy from the coverage of legislation such as this on any reasonable basis, just as with respect to any other physical condition.''
\end{quote}
\noindent --- Justice Potter Stewart, Opinion of the Court, 417 U.S. at 496--497 n. 20 (\href{https://supreme.justia.com/cases/federal/us/417/484/}{Justia}; \href{https://tile.loc.gov/storage-services/service/ll/usrep/usrep417/usrep417484/usrep417484.pdf}{Library of Congress U.S.\ Reports}).

\begin{quote}
``The lack of identity between the excluded disability and gender as such under this insurance program becomes clear upon the most cursory analysis. The program divides potential recipients into two groups---pregnant women and nonpregnant persons. While the first group is exclusively female, the second includes members of both sexes.''
\end{quote}
\noindent --- Justice Potter Stewart, Opinion of the Court, 417 U.S. at 496--497 n. 20 (\href{https://supreme.justia.com/cases/federal/us/417/484/}{Justia}).

\begin{quote}
``In my view, by singling out for less favorable treatment a gender-linked disability peculiar to women, the State has created a double standard for disability compensation: a limitation is imposed upon the disabilities for which women workers may recover, while men receive full compensation for all disabilities suffered, including those that affect only or primarily their sex, such as prostatectomies, circumcision, hemophilia, and gout. In effect, one set of rules is applied to females and another to males. Such dissimilar treatment of men and women, on the basis of physical characteristics inextricably linked to one sex, inevitably constitutes sex discrimination.''
\end{quote}
\noindent --- Justice William J. Brennan, Jr., dissenting, joined by Justices Douglas and Marshall, 417 U.S. at 501 (\href{https://supreme.justia.com/cases/federal/us/417/484/}{Justia}).

\end{document}
